\chapter{Mechanism Glossary}
\label{chap:glossary}

\section{Why this chapter exists}

The same words appear repeatedly in the runtime: task, shard, waker, readiness,
completion, submitter, snapshot, and scheduling group. This chapter is a short
field guide for those words.

\section{Runtime terms}

\begin{description}
    \item[Shard] A single ownership domain. In the standard-library services it
        is a worker thread with a mailbox. In \texttt{ShardedExecutor} it is a
        thread running one single-threaded executor.

    \item[Executor] A single-threaded async kernel. It owns tasks, timers,
        readiness registrations, optional \texttt{io\_uring} dispatcher state,
        and scheduler counters.

    \item[Task] The executor-owned wrapper around a pinned future. A task has
        an id, optional name, scheduling group id, lifecycle state, poll
        counters, timestamps, and a waker path back to its executor.

    \item[Root future] The future passed to \texttt{run\_until}. The executor
        drives it while also polling spawned tasks. When the root future
        completes, \texttt{run\_until} returns its output.

    \item[Spawner] A handle that submits new tasks to one executor. Cloned
        spawners keep the executor accepting work until they are dropped.

    \item[Submitter] A sharded capability for explicit cross-shard async
        submission. It owns the ability to place work on shard executors.

    \item[Shard mailbox] A bounded typed queue for transferring owned messages
        to one shard-local receiver. \texttt{ShardMailboxSet<M>} is the
        uniform form: one inbound mailbox per shard.

    \item[Work unit] A logical receiver role, such as an index assembler. A
        work unit has a name and an assigned shard. Multiple work units may be
        assigned to the same shard.

    \item[Work-unit mailbox] A named mailbox for non-uniform work assignment.
        The work-unit mailbox set routes by logical name and checks that the
        receiver is taken on the assigned shard.

    \item[Router] The key discriminator for message placement.
        \texttt{UniformShardRouter} maps a key to \texttt{ShardId};
        \texttt{WorkUnitRouter<N>} maps a key to a logical work-unit name.

    \item[Join handle] An awaitable handle for a spawned task result. It is
        also the cancellation handle: aborting it requests that the task be
        dropped if it has not completed. Dropping a join handle without
        awaiting it automatically aborts the underlying task.

    \item[Task scope] A structured owner for child tasks. It gives a group of
        children one cooperative stop source and one bounded shutdown path.
        Dropping a scope requests stop and aborts remaining children.
\end{description}

\section{Scheduling terms}

\begin{description}
    \item[Ready queue] A queue of tasks that should be polled. In the current
        executor, each scheduling group owns its own ready queue.

    \item[Waker] The object a future uses to ask the executor to poll it again.
        In \texttt{sitas}, waking a task requeues it on the same executor.

    \item[Poll] One call into a future. A future either returns
        \texttt{Ready(output)} or \texttt{Pending}. Long polls block the whole
        executor thread, so cooperative code must yield or await.

    \item[Scheduling group] An executor-local resource class with a name and
        relative shares. Tasks in a group are charged for observed poll time.

    \item[Virtual runtime] A weighted counter used to compare scheduling
        groups. More shares make the same poll duration charge less virtual
        runtime.

    \item[Ready-poll budget] The maximum number of ready tasks polled before
        the executor checks timers, readiness, and completions again.
\end{description}

\section{I/O terms}

\begin{description}
    \item[Readiness] The operating system reports that a descriptor may be
        ready for reading or writing. The runtime must still retry the
        operation and handle partial progress or \texttt{WouldBlock}.

    \item[Completion] The operating system reports that a submitted operation
        has finished. Linux \texttt{io\_uring} is completion-based.

    \item[Interest] A registered wish to be woken when a descriptor becomes
        readable or writable. Interests are task/waker bookkeeping, not file
        descriptor ownership.

    \item[Dispatcher] The executor-local \texttt{io\_uring} state machine that
        owns submission tracking, completion buffering, waker registration, and
        abandoned operation cleanup.

    \item[Abandoned operation] A completion operation whose future was dropped
        before the kernel completion was observed. The dispatcher must keep
        enough state to dispose of buffers safely later.
\end{description}

\section{Observability terms}

\begin{description}
    \item[Running statistics] An online accumulator using Welford's algorithm
        that computes count, mean, variance, standard deviation, coefficient
        of variation, min, max, and total from streaming \texttt{f64} samples
        without storing the values. Two accumulators can be merged. Duration
        convenience methods accept \texttt{Duration} and \texttt{Instant} pairs.

    \item[Snapshot] An owned point-in-time copy of runtime state. A snapshot is
        not live and does not keep the runtime alive.

    \item[Coefficient of variation (CV)] Standard deviation divided by mean,
        expressed as a percentage. Unit-free, useful for comparing dispersion
        across metrics with different scales.

    \item[Runtime identity] A private identity used to reject handles from a
        different runtime, such as a sharded scheduling group created by
        another \texttt{ShardedExecutor}.
\end{description}

\section{Ownership terms}

\begin{description}
    \item[Owned value] A value that can safely cross a boundary because no
        borrowed reference into local runtime or service state escapes with it.

    \item[Shard-local value] One owned value per shard, accessed only on the
        owning shard and only through synchronous closures.
        \texttt{with\_current\_result} skips the \texttt{ShardId} parameter.
        \texttt{get\_current} returns a clone.

    \item[Stoppable shard-local workers] A set of shard-local workers sharing
        a cooperative stop token. Dropping the handle requests stop and aborts
        remaining workers---the explicit \texttt{stop\_and\_join} path is
        preferred for clean shutdown.
\end{description}

\section{Mechanism maps}

\subsection{Wake path}

\begin{center}
\begin{tikzpicture}[
    mechbox/.style={draw=sitasMuted, fill=sitasCodeBackground,
        minimum width=82mm, minimum height=8mm, text width=78mm,
        align=center, font=\footnotesize\ttfamily, rounded corners=1.5pt},
    arr/.style={->, >=stealth, thick, color=sitasAccent}
]
\node[mechbox] (waker) at (0, 0) {Future stores Waker};
\node[mechbox] (event) at (0, -1.1) {external event or explicit wake};
\node[mechbox] (queued) at (0, -2.2) {Task waker marks task queued};
\node[mechbox] (group) at (0, -3.3) {Scheduler pushes task into its scheduling group queue};
\node[mechbox] (poll) at (0, -4.4) {Executor polls ready work later};

\draw[arr] (waker) -- (event);
\draw[arr] (event) -- (queued);
\draw[arr] (queued) -- (group);
\draw[arr] (group) -- (poll);
\end{tikzpicture}
\end{center}

\subsection{Timer path}

\begin{center}
\begin{tikzpicture}[
    mechbox/.style={draw=sitasMuted, fill=sitasCodeBackground,
        minimum width=82mm, minimum height=8mm, text width=78mm,
        align=center, font=\footnotesize\ttfamily, rounded corners=1.5pt},
    arr/.style={->, >=stealth, thick, color=sitasAccent}
]
\node[mechbox] (sleep) at (0, 0) {sleep(duration)};
\node[mechbox] (register) at (0, -1.1) {register deadline and waker};
\node[mechbox] (wait) at (0, -2.2) {executor waits until earliest deadline or another event};
\node[mechbox] (expired) at (0, -3.3) {expired timer wakes task};
\node[mechbox] (requeued) at (0, -4.4) {task is requeued and polled};

\draw[arr] (sleep) -- (register);
\draw[arr] (register) -- (wait);
\draw[arr] (wait) -- (expired);
\draw[arr] (expired) -- (requeued);
\end{tikzpicture}
\end{center}

\subsection{Readiness path}

\begin{center}
\begin{tikzpicture}[
    mechbox/.style={draw=sitasMuted, fill=sitasCodeBackground,
        minimum width=82mm, minimum height=8mm, text width=78mm,
        align=center, font=\footnotesize\ttfamily, rounded corners=1.5pt},
    arr/.style={->, >=stealth, thick, color=sitasAccent}
]
\node[mechbox] (try) at (0, 0) {read/write future tries non-blocking operation};
\node[mechbox] (wouldblock) at (0, -1.1) {WouldBlock};
\node[mechbox] (interest) at (0, -2.2) {register fd interest and task waker};
\node[mechbox] (ready) at (0, -3.3) {reactor reports readiness};
\node[mechbox] (retry) at (0, -4.4) {task retries operation};

\draw[arr] (try) -- (wouldblock);
\draw[arr] (wouldblock) -- (interest);
\draw[arr] (interest) -- (ready);
\draw[arr] (ready) -- (retry);
\end{tikzpicture}
\end{center}

\subsection{Completion path}

\begin{center}
\begin{tikzpicture}[
    mechbox/.style={draw=sitasMuted, fill=sitasCodeBackground,
        minimum width=82mm, minimum height=8mm, text width=78mm,
        align=center, font=\footnotesize\ttfamily, rounded corners=1.5pt},
    arr/.style={->, >=stealth, thick, color=sitasAccent}
]
\node[mechbox] (submit) at (0, 0) {io\_uring future submits operation with owned buffer};
\node[mechbox] (track) at (0, -1.1) {dispatcher tracks operation id, buffer, and waker};
\node[mechbox] (kernel) at (0, -2.2) {kernel completes operation};
\node[mechbox] (dispatch) at (0, -3.3) {dispatcher buffers completion or wakes task};
\node[mechbox] (take) at (0, -4.4) {future takes result and buffer};

\draw[arr] (submit) -- (track);
\draw[arr] (track) -- (kernel);
\draw[arr] (kernel) -- (dispatch);
\draw[arr] (dispatch) -- (take);
\end{tikzpicture}
\end{center}

\section{Common confusions}

\begin{description}
    \item[Readiness is not completion] Readiness says retrying may work.
        Completion says submitted work has finished.

    \item[Shared-nothing is not synchronization-free] The runtime uses locks
        and atomics internally. The application-state model remains
        shared-nothing.

    \item[A scheduling group is not a thread] It is a queue and accounting
        class inside one executor. On the sharded runtime, a sharded group is
        one matching local group per shard.

    \item[A snapshot is not a subscription] It is an owned copy. To observe
        progress over time, take repeated snapshots and compare them.

    \item[CPU placement is not load balancing] Placement asks the operating
        system where shard threads should run. It does not redistribute work
        between shards.
\end{description}
